\chapter{Abordagem proposta}
	\label{chap:propApproach}

	\section{A Base de sinais}
   		\label{sec:baseDeSinais}

		\par Os experimentos deste trabalho utilizam a base pública de sinais descrita em \cite{10.1117/12.2255697}, que reúne registros simultâneos de voz e eletroencefalograma (EEG) de 15 voluntários falantes nativos de espanhol. Cada participante produziu onze classes de enunciados: cinco vogais (/a/, /e/, /i/, /o/, /u/) e seis comandos direcionais (\textit{arriba}, \textit{abajo}, \textit{adelante}, \textit{atras}, \textit{derecha}, \textit{izquierda}), em duas modalidades --- fala fonatória e fala imaginada.

		\par Os sinais de voz foram gravados a 44100\,Hz com canal único. Os sinais de EEG foram capturados com seis eletrodos posicionados em F3, F4, C3, C4, P3 e P4, conforme o sistema internacional 10-20, com taxa de amostragem de 1024\,Hz. Cada trial corresponde a uma janela de 4 segundos da fase de imaginar/pronunciar, resultando em 4096 amostras por canal de EEG e 176400 amostras de áudio por trial.

		\par Os dados são organizados em arquivos MATLAB (\texttt{.mat}) por sujeito, com dois arquivos por participante: \texttt{S\#\#\_EEG.mat} e \texttt{S\#\#\_Audio.mat}. Cada arquivo contém uma matriz cujas linhas representam trials individuais. No arquivo de EEG, cada linha possui 24576 colunas de sinal (6 canais $\times$ 4096 amostras, concatenados) seguidas de 3 colunas de metadados: modalidade (imaginada ou fonatória), código do estímulo e indicador de artefato de piscada ocular. No arquivo de áudio, cada linha contém 176400 amostras seguidas de 2 colunas: código do estímulo e índice da linha correspondente na matriz de EEG, garantindo o alinhamento temporal entre as duas modalidades. A \autoref{fig:baseDeSinais} ilustra essa organização.

		\begin{figure}[ht]
			\centering
			\caption[Organização da base de sinais]{Estrutura dos arquivos da base de sinais \cite{10.1117/12.2255697}}
			\scalebox{0.82}{
				\begin{tikzpicture}[
    rootfolder/.style={rectangle, rounded corners=3pt, draw, fill=orange, text=black,
        minimum width=4.4cm, minimum height=0.65cm,
        font=\footnotesize\sffamily\bfseries, align=center},
    subfolder/.style={rectangle, rounded corners=3pt, draw, fill=orange!70, text=black,
        minimum width=3.4cm, minimum height=0.60cm,
        font=\scriptsize\sffamily\bfseries, align=center},
    txtfile/.style={rectangle, draw, fill=gray!30, text=black,
        minimum width=3.4cm, minimum height=0.55cm,
        font=\scriptsize\sffamily, align=center},
    matfile/.style={rectangle, rounded corners=2pt, draw, fill=green!60, text=black,
        minimum width=3.2cm, minimum height=0.55cm,
        font=\scriptsize\sffamily\bfseries, align=center},
    sigblock/.style={rectangle, draw, fill=green!60, text=black,
        minimum height=0.65cm,
        font=\scriptsize\sffamily\bfseries, align=center, inner sep=4pt},
    metablock/.style={rectangle, draw, fill=yellow, text=black,
        minimum height=0.65cm,
        font=\tiny\sffamily\bfseries, align=center, inner sep=2pt},
    refblock/.style={rectangle, draw, fill=orange, text=black,
        minimum height=0.65cm,
        font=\tiny\sffamily\bfseries, align=center, inner sep=2pt},
    arr/.style={->, >=stealth, draw=gray, line width=1.0mm},
    darr/.style={->, >=stealth, dashed, draw=gray, line width=1.0mm},
    treeline/.style={draw=gray!70, line width=0.5mm},
    lbl/.style={font=\tiny\sffamily\bfseries, text=black},
]

%% ===== FOLDER TREE =====
\node[rootfolder] (root) at (2.2, 0) {Imagined Speech Database};

%% Root-level items (indent x=0.5)
\node[txtfile, anchor=west]   (proto) at (0.5, -0.80) {Protocol.txt};
\node[txtfile, anchor=west]   (subjs) at (0.5, -1.50) {Subjects.txt};
\node[subfolder, anchor=west] (s01d)  at (0.5, -2.20) {S01/};

%% S01 contents (indent x=1.1)
\node[txtfile, minimum width=3.2cm, anchor=west] (s01txt) at (1.1, -3.00) {S01.txt};
\node[matfile, anchor=west] (s01eeg) at (1.1, -3.70) {S01\_EEG.mat};
\node[matfile, anchor=west] (s01aud) at (1.1, -4.40) {S01\_Audio.mat};

%% Remaining subjects
\node[subfolder, anchor=west] (s02d) at (0.5, -5.10) {S02/};
\node[lbl, anchor=west] at (0.85, -5.75) {$\vdots$};
\node[subfolder, anchor=west] (s15d) at (0.5, -6.30) {S15/};

%% Tree lines: root trunk
\draw[treeline] (0.20,-0.35) -- (0.20,-6.30);
\foreach \cy in {-0.80,-1.50,-2.20,-5.10,-6.30} {
    \draw[treeline] (0.20,\cy) -- (0.48,\cy);
}
%% S01 subtree trunk
\draw[treeline] (0.80,-2.50) -- (0.80,-4.40);
\foreach \cy in {-3.00,-3.70,-4.40} {
    \draw[treeline] (0.80,\cy) -- (1.08,\cy);
}

%% ===== MATRIX ROW DIAGRAMS (x starts at 4.7) =====

%% EEG matrix row
\node[sigblock, anchor=west, minimum width=4.8cm] (eegsig)
    at (4.7, -3.20)
    {Sinal EEG \\ \tiny $6 \times 4096 = 24576$ col.};
\node[lbl, below=0.03cm of eegsig.south, anchor=north]
    {\tiny F3\;F4\;C3\;C4\;P3\;P4};
\node[metablock, anchor=west, minimum width=0.80cm] (eegmode) at (eegsig.east) {modo};
\node[metablock, anchor=west, minimum width=0.80cm] (eegstim) at (eegmode.east) {est.};
\node[metablock, anchor=west, minimum width=1.00cm] (eegblink) at (eegstim.east) {pisc.};
\node[lbl, above=0.04cm of eegsig.north west, anchor=south west]
    {linha\,$=$\,1 trial (4\,s)};

%% Audio matrix row
\node[sigblock, anchor=west, minimum width=4.8cm] (audiosig)
    at (4.7, -5.20)
    {Sinal de áudio \\ \tiny $176400$ colunas};
\node[metablock, anchor=west, minimum width=0.80cm] (audiostim) at (audiosig.east) {est.};
\node[refblock, anchor=west, minimum width=2.50cm] (audioidx) at (audiostim.east) {índice EEG};
\node[lbl, above=0.04cm of audiosig.north west, anchor=south west]
    {linha\,$=$\,1 trial (4\,s)};

%% Arrows: .mat files → matrix rows
\draw[arr] (s01eeg.east) to[out=0,in=180] (eegsig.west);
\draw[arr] (s01aud.east) to[out=0,in=180] (audiosig.west);

%% Cross-reference arrow: audio index → EEG row
\draw[darr] (audioidx.east) -- ++(0.35,0) |- (eegblink.east);
\node[lbl, rotate=90, anchor=south]
    at ([xshift=0.48cm, yshift=1.0cm] audioidx.east)
    {ref.\ cruzada};

\end{tikzpicture}

			}
			\legend{Fonte: Elaborado pelo autor, 2025.}
			\label{fig:baseDeSinais}
		\end{figure}

	\section{Estrutura da estratégia proposta}
		\label{sec:estruturaDaEstrategiaProposta}
		
		\par Primeiramente, os sinais são decompostos pela \ac{DTWPT}; em seguida, realizam-se os agrupamentos energéticos em três esquemas de particionamento espectral --- bandas lineares, Mel e Bark --- gerando, na Categoria~1, os vetores de \textit{Linear-band energies}, \textit{Mel-band energies} e \textit{Bark-band energies}, respectivamente. Para a Categoria~2, aplica-se adicionalmente o logaritmo e a \ac{DCT} sobre cada um desses vetores, produzindo os coeficientes \ac{LFCC}, \ac{MFCC} e \ac{BFCC}. Por fim, conduz-se a análise paraconsistente de características a fim de selecionar as melhores combinações entre agrupamento energético e \textit{wavelet}, visando gerar características as mais disjuntas possíveis.

		\par Em segundo lugar, tanto os dados de voz quanto os de EEG são submetidos a \textit{autoencoders} para extração de características --- a Categoria~3 da taxonomia introduzida na \autoref{sec:conceitos}. As três categorias são, portanto, as três famílias de extratores comparadas nesta tese.

		\par As características extraídas, tanto automatizadas quanto as manualmente obtidas, servem de entrada para uma rede neural residual profunda de pulsos (\ac{DSNN}), que determina as classes de cada um dos sinais. A arquitetura dessa rede e o protocolo de avaliação são detalhados na \autoref{sec:fase01}.

		\par Como o tema central deste trabalho é a autenticação biométrica combinando fala fonada e fala imaginada, cada configuração é avaliada em três modos, para medir a contribuição individual e conjunta das duas fontes:
		\begin{itemize}
			\item \textbf{apenas fala fonada}: somente o sinal de voz (áudio) alimenta o classificador;
			\item \textbf{apenas fala imaginada}: somente o sinal de EEG alimenta o classificador;
			\item \textbf{fonada + imaginada (fundida)}: as duas fontes são combinadas antes da classificação.
		\end{itemize}

		\par No modo fundido, o \emph{instante} em que as duas fontes se combinam é, ele próprio, um eixo experimental, pois desloca a fronteira entre o que é modelado separadamente e o que é modelado em conjunto:
		\begin{itemize}
			\item \textbf{fusão precoce} (\textit{early fusion}): os sinais brutos de voz e EEG são concatenados \emph{antes} da extração de características, de modo que um único extrator opera sobre o sinal conjunto e pode capturar interações entre as modalidades. Em contrapartida, impõe uma taxa de amostragem comum às duas metades: adota-se a taxa da voz (44100\,Hz), dominante em número de amostras, o que exige reamostrar o EEG de 1024\,Hz para 44100\,Hz --- um fator de aproximadamente $43\times$. Essa reamostragem não cria informação: ela apenas interpola o sinal de EEG para que as duas metades possam ocupar um mesmo vetor, e o custo é que a metade de EEG do sinal concatenado passa a ser majoritariamente composta por amostras interpoladas;
			\item \textbf{fusão tardia} (\textit{late fusion}): cada sinal é submetido \emph{independentemente} à extração de características, e apenas os vetores resultantes são concatenados. Preserva a taxa nativa e o extrator próprio de cada fonte, ao custo de nunca observar as duas conjuntamente.
		\end{itemize}
		\par Ambas as estratégias são comparadas, uma vez que não há garantia a priori de qual delas produz características mais separáveis segundo a engenharia paraconsistente.

		\par Como o modo fundido se desdobra nessas duas variantes, os \emph{três} modos acima correspondem, no eixo experimental da Fase~01 (\autoref{sec:fase01}), a \emph{quatro} fontes distintas: apenas voz, apenas EEG, fundida-precoce e fundida-tardia.

		\par Quanto à extração automatizada, comparam-se dois tipos de \textit{autoencoder} --- um de pulsos (\ac{SNN}-AE) e um convencional (\ac{ANN}-AE) --- e ambos são confrontados com a extração manual (\textit{handcrafted}), usando a engenharia paraconsistente de características como critério de comparação (\autoref{sec:conceitos}).

		\par Uma visão geral do processo é mostrada na \autoref{fig:estruturaDaEstrategiaProposta}.
		
		\begin{figure}[h]
			\centering
			\caption[Estratégia proposta]{Fluxograma da estratégia proposta}
			\scalebox{0.75}	{
				\input{images/strategicStructure.tex}
			}
			\legend{Fonte: Elaborado pelo autor, 2024.}
			\label{fig:estruturaDaEstrategiaProposta}
		\end{figure}

		\subsection{Protocolo experimental em duas fases}
			\label{sec:protocoloDuasFases}

			\par O experimento é dividido em duas fases independentes, executadas em sequência. A \textbf{Fase 00} constrói o vetor de características e escolhe, por sinal, o melhor extrator segundo a engenharia paraconsistente. A \textbf{Fase 01} recebe apenas o vetor vencedor e treina o classificador para medir o desempenho de autenticação. A separação é controlada por um único parâmetro de configuração, \texttt{classifier.enabled}: quando falso, a execução para após o ranking (Fase 00); quando verdadeiro, o classificador é treinado (Fase 01). A \autoref{fig:e05PipelineTwoPhase} mostra o encadeamento completo.

			\begin{figure}[h]
				\centering
				\caption[Pipeline em duas fases]{Encadeamento das duas fases e dos \textit{scripts} de transição}
				\scalebox{0.9}{\begin{tikzpicture}[
    startstop/.style={
        rectangle, rounded corners=4pt, draw, fill=flowStartEndColor,
        text=black, minimum width=6.0cm, minimum height=0.7cm,
        font=\footnotesize\sffamily\bfseries, align=center
    },
    proc/.style={
        rectangle, rounded corners, draw, fill=flowProcessColor,
        minimum width=6.0cm, minimum height=0.8cm,
        font=\footnotesize\sffamily\bfseries, align=center
    },
    script/.style={
        rectangle, draw, fill=cbBlue!25,
        minimum width=6.0cm, minimum height=0.8cm,
        font=\footnotesize\sffamily\bfseries, align=center
    },
    arr/.style={->, >=stealth, draw=flowArrowColor, line width=1.0mm},
    node distance=0.55cm
]

\node[startstop] (start) {Base de sinais (voz + EEG pareados)};
\node[proc, below=of start] (p00)
    {\textbf{Fase 00} --- varrer $23$ wavelets $\times$ $3$ escalas $\times$ $2$ categorias $\times$ $2$ sinais\\ ($+$ 24 autoencoders compactos --- SNN-AE $\times$3 codificações e ANN-AE): $300$ execuções, \texttt{classifier.enabled=false}};
\node[script, below=of p00] (rank)
    {\texttt{e05\_phase00\_rank.py}\\ menor $D_{\text{verdade}}$ por sinal $\rightarrow$ \texttt{winners.json}};
\node[script, below=of rank] (apply)
    {\texttt{e05\_apply\_winner.py}\\ injeta o vencedor nos $32$ perfis da Fase 01};
\node[proc, below=of apply] (p01)
    {\textbf{Fase 01} --- DSNN sobre $4$ fontes $\times$ $2$ textos $\times$ $2$ CV\\ $\times$ $2$ padronização: $32$ execuções, \texttt{classifier.enabled=true}};
\node[startstop, below=of p01] (stop) {Métricas de verificação (EER, AUC)};

\draw[arr] (start) -- (p00);
\draw[arr] (p00)   -- (rank);
\draw[arr] (rank)  -- (apply);
\draw[arr] (apply) -- (p01);
\draw[arr] (p01)   -- (stop);
\end{tikzpicture}
}
				\legend{Fonte: Elaborado pelo autor, 2026.}
				\label{fig:e05PipelineTwoPhase}
			\end{figure}

			\subsubsection{Fase 00 --- construção e seleção do vetor de características}

				\par A Fase 00 é aplicada separadamente à voz e ao EEG. Para o método manual, varrem-se duas escolhas: a \textbf{wavelet-mãe} da DTWPT (23 opções: \texttt{haar} e \texttt{daub4} a \texttt{daub46}) e a \textbf{escala} de agrupamento espectral (\textit{bark}, \textit{mel}, \textit{lfcc}). Além disso, escolhe-se a \textbf{categoria} do descritor: Categoria~1 usa as energias de banda diretamente; Categoria~2 aplica logaritmo e DCT sobre essas energias, gerando os coeficientes cepstrais LFCC, MFCC ou BFCC (para escala \textit{lfcc}, \textit{mel} ou \textit{bark}). Para a \textbf{voz}, isso dá $23 \times 3 \times 2 = 138$ combinações manuais. Para o \textbf{EEG}, porém, o eixo da escala não se aplica: \textit{bark} e \textit{mel} são escalas cocleares --- modelam a resolução em frequência da \emph{audição} humana --- e não há base fisiológica para aplicá-las a um sinal que não é som. Essa exclusão é também empiricamente inócua: sobre o agrupamento espectral usado aqui, as três escalas produzem para o EEG exatamente o mesmo particionamento (\autoref{sec:escalaEeg}), de modo que as três variantes eram numericamente idênticas. O EEG usa portanto apenas a escala linear (\textit{lfcc}), dando $23 \times 1 \times 2 = 46$ combinações manuais. A ambos os sinais somam-se $12$ \textit{autoencoders} compactos, sem wavelet ou escala: $9$ SNN-AE ($3$ tamanhos $\times$ $3$ codificações temporais) e $3$ ANN-AE ($3$ tamanhos). Ambas as famílias consomem um vetor achatado de dimensão fixa (o sinal é reamostrado por \textit{pooling} para $256$ amostras, normalizado min-max para $[0,1]$), usam uma única camada oculta (rasa, adequada a dispositivos de baixo poder) e razão de compressão $2\!:\!1$ entre a camada oculta e o gargalo latente (dimensões latentes $8$, $16$ e $32$, tamanhos \texttt{tiny}/\texttt{small}/\texttt{base}). O latente é o próprio vetor de características, então um gargalo menor gera vetor menor, barateando o classificador a jusante --- princípio da literatura de autoencoders leves para borda.

\textbf{SNN-AE} (codificador \texttt{Linear}$\to$\texttt{LIF}, decodificador \texttt{Linear}$\to$\texttt{LIF-integrador}) só carrega informação através de pulsos: um único vetor analógico apresentado em um só passo (equivalente a $T=1$) deixa o neurônio LIF quase sempre abaixo do limiar de disparo, colapsando o latente a zero. Por isso cada amostra normalizada é expandida em $T$ (padrão $16$) quadros de pulso via três codificações temporais comparadas nesta tese --- $18 = 3~\text{tamanhos} \times 3~\text{codificações} \times 2~\text{sinais}$:
\begin{itemize}
	\item \texttt{poisson} --- cada componente dispara com probabilidade de Bernoulli igual ao seu valor normalizado (codificação por taxa);
	\item \texttt{latency} --- componentes de maior valor disparam mais cedo ($t_{\text{disparo}} = \text{round}((1-v)(T-1))$, mesmo esquema do Experimento 04);
	\item \texttt{direct} --- passagem analógica direta, sem pulsos (linha de base não-temporal; $T=1$).
\end{itemize}
O estado de membrana é \emph{reiniciado uma única vez por amostra} e depois \emph{integra ao longo dos $T$ quadros} (sem reset entre quadros) --- é essa integração temporal que permite a uma corrente sub-limiar por passo acumular até disparar. A característica de cada amostra é a \emph{taxa de disparo média do latente sobre os $T$ quadros}. O limiar de disparo do codificador (\texttt{voltage\_threshold}) é reduzido para $0{,}2$ (padrão LIF $=1{,}0$) nas variantes \texttt{poisson}/\texttt{latency}, já que quadros de pulso normalizados têm amplitude baixa; \texttt{direct} mantém o limiar padrão $1{,}0$, reproduzindo intencionalmente o regime não-temporal, que serve de linha de base para as duas codificações temporais.

\paragraph{Codificação temporal e função de perda do treinamento.} Embora as três codificações produzam padrões de pulso distintos na entrada do codificador, o SNN-AE é treinado com a \textbf{mesma} função de perda de reconstrução --- Erro Médio Quadrático (EQM) --- nas três variantes (\texttt{poisson}, \texttt{latency}, \texttt{direct}), pois o alvo da reconstrução é sempre o sinal analógico original de entrada, e não um trem de pulsos: a codificação temporal afeta apenas como a informação flui pelo codificador, não o que o decodificador precisa reproduzir. É importante notar que o \textit{framework} desenvolvido para esta tese também implementa duas funções de perda casadas a codificações de pulso --- \texttt{SpikeCountLoss} (EQM sobre a contagem total de disparos, adequada quando o próprio alvo supervisionado é uma taxa de disparo) e \texttt{SpikeTimeLoss} (EQM sobre o instante do primeiro disparo, adequada quando o alvo é um único evento temporal) --- seguindo a convenção da codificação temporal usada por \cite{comsa2021spiking} e a discussão de funções de perda sob medida para predição de pulsos de \cite{manna2024time}. Usar a função casada errada quebra o treino: \texttt{SpikeTimeLoss} sobre um alvo de taxa enxerga apenas o primeiro disparo e ignora a contagem; \texttt{SpikeCountLoss} sobre um alvo de latência trata a ausência de disparo como contagem zero, gerando gradiente incorreto para disparos tardios. Essa correspondência codificação--perda é, no entanto, uma capacidade geral do \textit{framework} (\texttt{include/layers/losses/SpikeCountLoss.hpp}, \texttt{SpikeTimeLoss.hpp}), exercitada apenas nos testes unitários da biblioteca --- o SNN-AE da Fase~00 não a utiliza, precisamente porque seu alvo de treino nunca é um trem de pulsos, e sim o sinal contínuo original.

% TODO (figura pendente): ilustrar as 3 codificações temporais
% (poisson/latency/direct) aplicadas a uma imagem de teste, 3 sub-figuras lado a
% lado, cada uma mostrando os T quadros de pulso empilhados. Gerar com script
% Python (matplotlib) reaproveitando spike_frame() de E05FeatureExtraction.cpp.
% O float foi REMOVIDO (e não deixado vazio) porque uma figura sem imagem é
% impressa como um quadro em branco numerado e referenciado no texto.

\paragraph{Comparação entre as codificações temporais.} As três codificações foram executadas sobre a grade da Fase~00 (18 perfis SNN-AE --- 3 codificações $\times$ 3 tamanhos $\times$ 2 sinais, 3 repetições cada). Uma versão anterior deste texto reportava valores de $\alpha$ e $D_{\text{verdade}}$ medidos sob o defeito de escala da taxa de aprendizado descrito em \autoref{sec:snnLrScale} --- a redução destinada apenas aos parâmetros biofísicos era aplicada a \emph{todos} os parâmetros, de modo que cada peso treinou a uma taxa efetiva dez vezes menor que a declarada no perfil. Aqueles números foram \textbf{retirados}, e não corrigidos: descreviam um experimento diferente do que os perfis declaravam. Corrigido o defeito, os resultados armazenados foram descartados e a Fase~00 foi integralmente reexecutada; os valores atuais estão nas Tabelas~\ref{tab:phase00ae_eeg} e~\ref{tab:phase00ae_voice}.

A reexecução confirma a ordem \texttt{direct} $\succ$ \texttt{latency} $\succ$ \texttt{poisson} em ambos os sinais, e de forma limpa: as três codificações formam blocos contíguos no ranking, com a \texttt{poisson} ocupando as três últimas posições e exibindo $\alpha$ praticamente nulo ($\alpha \le 0{,}003$ nos seis perfis). \textbf{Isso não estabelece um resultado negativo contra a codificação temporal}, e a razão é o argumento desenvolvido a seguir: essa é exatamente a ordem que o piso de ruído de cada codificação prevê, \emph{antes de qualquer aprendizado}. O que a medição confirma é a previsão do mecanismo, não uma propriedade informacional dos códigos temporais.

O \emph{mecanismo} que sustenta essa leitura é teórico, e prevê a ordem observada independentemente de qualquer execução --- é por isso que a concordância entre previsão e medição, por si só, não decide a questão. Cada codificação injeta, \textbf{antes de qualquer aprendizado}, uma quantidade de ruído que se calcula em fechado. A \texttt{poisson} sorteia um Bernoulli independente a cada passo de tempo, de modo que a média sobre $T$ quadros tem variância $v(1-v)/T$; para $v = 0{,}5$ (pior caso) e $T = 16$, isso dá desvio-padrão

\begin{equation}
	\label{eq:ruidoPoisson}
	\sigma_{\text{poisson}} = \sqrt{\frac{0{,}25}{16}} = 0{,}125 \: ,
\end{equation}

isto é, \textbf{$12{,}5\%$ da faixa total $[0,1]$ da característica}, somados a cada amostra de forma independente e aleatória --- e presentes ainda que o codificador fosse perfeito, pois nascem na entrada. A \texttt{latency}, ao contrário, é determinística: o mesmo $v$ sempre produz o mesmo instante de disparo, e seu erro é de quantização, limitado a meio degrau, $\tfrac{1}{2}\cdot\tfrac{1}{T-1} \approx 0{,}033$ --- quase quatro vezes menor e, sobretudo, idêntico entre amostras de mesmo valor. A \texttt{direct} não injeta ruído algum.

Esse escalonamento --- ruído nulo, pequeno-determinístico, grande-aleatório --- incide diretamente sobre o que $\alpha$ mede. Como $\alpha$ é definido a partir da \emph{amplitude} (mínimo--máximo) intraclasse (\autoref{sec:conceitos}), ele é maximamente sensível a valores extremos: basta que a estocasticidade produza poucas amostras extremas para que a amplitude de uma classe cubra todo o intervalo normalizado e $\alpha \to 0$. O piso de ruído de cada codificação mapeia-se, portanto, monotonicamente em $\alpha$, e prevê a ordem \texttt{direct} $\succ$ \texttt{latency} $\succ$ \texttt{poisson} como \textbf{artefato conjunto da métrica e do piso estatístico do codificador}, não como evidência de que códigos temporais carreguem menos informação de locutor.

Daí a retificação metodológica: a afirmação de que a codificação temporal constitui um resultado negativo \textbf{não está estabelecida}, e não por mera incerteza --- $\alpha$ é \emph{estruturalmente incapaz} de distinguir ``o codificador não aprendeu'' de ``esta codificação tem um piso estatístico irredutível em $T=16$''. Como esse piso decresce com $1/T$ (\autoref{eq:ruidoPoisson}), a questão é empiricamente testável: em $T=64$ a variância cai por um fator de quatro, e $\sigma_{\text{poisson}}$ vai a $0{,}0625$. Esse teste ainda não foi executado. Registre-se também que este problema é \textbf{independente} da vulnerabilidade de $D_{1,0}$ tratada em \autoref{sec:dPenalizado}: aquela mostra que o critério \emph{recompensa} indevidamente variância nula; esta, que ele pode \emph{punir} indevidamente uma variância real que não decorre de pior aprendizado. Corrigir a primeira não corrige a segunda.

Cada combinação gera um vetor de características, avaliado pela distância $D_{\text{verdade}}$ --- a distância ao vértice \textit{Verdade} no plano paraconsistente (\autoref{sec:conceitos}); quanto menor, mais separáveis são as classes antes de qualquer classificador. O vencedor do sinal, porém, \textbf{não} é a combinação de menor $D_{\text{verdade}}$, e sim a de menor $D_{\text{penalizado}} = D_{\text{verdade}} + \lambda|G_2|$, com $\lambda = 2-\sqrt{2}$ (\autoref{sec:dPenalizado}): $D_{\text{verdade}}$ isolado é vulnerável a extratores degenerados, que produzem uma saída praticamente constante e assim obtêm uma pontuação artificialmente boa sem carregar qualquer informação de classe. O termo $\lambda|G_2|$ penaliza exatamente essa degenerescência. A escolha altera o resultado na prática, e não apenas em princípio --- as consequências sobre o vencedor de cada sinal são discutidas no \autoref{chap:testsResults}. O vencedor é escolhido entre \emph{todas} as combinações do sinal, manuais e \textit{autoencoders} conjuntamente. Nenhum classificador é treinado nesta fase. O \autoref{lst:e05Phase00} descreve o procedimento e a \autoref{fig:e05Phase00Flow} o resume.

				\begin{lstlisting}[language=C, caption={Fase 00 --- varredura de extratores e seleção do vencedor por sinal via engenharia paraconsistente. Executada uma vez para a voz e uma vez para o EEG.}, label={lst:e05Phase00}]
/* WAVELETS: 23 wavelets-mae (haar, daub4, ..., daub46)          */
/* ESCALAS : 3 escalas espectrais (bark, mel, lfcc)              */
/* sinal   : vetor de amostras de um sinal (voz OU EEG)          */
/* Retorna a melhor combinacao (menor D_verdade).                */

Combinacao seleciona_vencedor(double sinal[], int n)
{
    Combinacao vencedor;
    vencedor.d_verdade = INFINITO;      /* comeca no pior valor  */

    /* 1) parte handcrafted: varre wavelet X escala             */
    for (int w = 0; w < 23; w++) {
        for (int e = 0; e < 3; e++) {
            double vetor[MAX_CARACT];
            int m = extrai_handcrafted(sinal, n, WAVELETS[w],
                                       ESCALAS[e], vetor);

            /* alpha, beta -> ponto no plano paraconsistente     */
            double d = calcula_d_verdade(vetor, m);

            if (d < vencedor.d_verdade) {   /* mais separavel    */
                vencedor.d_verdade = d;
                vencedor.estrategia = HANDCRAFTED;
                vencedor.wavelet = WAVELETS[w];
                vencedor.escala  = ESCALAS[e];
            }
        }
    }

    /* 2) parte automatizada: um autoencoder (sem wavelet/escala) */
    double vetor_ae[MAX_CARACT];
    int mae = extrai_autoencoder(sinal, n, vetor_ae);
    double d_ae = calcula_d_verdade(vetor_ae, mae);
    if (d_ae < vencedor.d_verdade) {
        vencedor.d_verdade = d_ae;
        vencedor.estrategia = AUTOENCODER;
    }

    return vencedor;                    /* NAO treina classificador */
}
\end{lstlisting}


				\begin{figure}[h]
					\centering
					\caption[Fluxograma da Fase 00]{Seleção do vetor vencedor de um sinal na Fase 00}
					\scalebox{0.82}{\begin{tikzpicture}[
    startstop/.style={
        rectangle, rounded corners=4pt, draw, fill=orange,
        text=black, minimum width=3.4cm, minimum height=0.7cm,
        font=\footnotesize\sffamily\bfseries, align=center
    },
    proc/.style={
        rectangle, rounded corners, draw, fill=green,
        minimum width=5.2cm, minimum height=0.7cm,
        font=\footnotesize\sffamily\bfseries, align=center
    },
    dec/.style={
        diamond, aspect=2.2, draw, fill=yellow!70,
        minimum width=3.0cm, inner sep=1pt,
        font=\scriptsize\sffamily\bfseries, align=center
    },
    arr/.style={->, >=stealth, draw=gray, line width=1.0mm},
    lbl/.style={font=\scriptsize\sffamily\bfseries, text=black},
    node distance=0.5cm and 1.4cm
]

\node[startstop] (start) {Sinal (voz ou EEG)};
\node[dec, below=of start] (loop)
    {resta par\\ (wavelet, escala)?};
\node[proc, below=1.1cm of loop] (extract)
    {DTWPT com a wavelet $\rightarrow$ agrupar na escala\\ $\rightarrow$ descritores $\rightarrow$ vetor de características};
\node[proc, below=of extract] (score)
    {engenharia paraconsistente:\\ calcular $\alpha$, $\beta$, $D_{\text{verdade}}$};
\node[proc, right=of loop] (pick)
    {vencedor $=$ combinação\\ de menor $D_{\text{verdade}}$};
\node[startstop, right=of pick] (stop) {Vetor vencedor do sinal};

\draw[arr] (start)   -- (loop);
\draw[arr] (loop)    -- node[lbl, right] {sim} (extract);
\draw[arr] (extract) -- (score);
\draw[arr] (score.west) -- ++(-1.4cm,0) |- (loop.west);
\draw[arr] (loop)    -- node[lbl, above] {não} (pick);
\draw[arr] (pick)    -- (stop);
\end{tikzpicture}
}
					\legend{Fonte: Elaborado pelo autor, 2026.}
					\label{fig:e05Phase00Flow}
				\end{figure}

				\par Considerando os dois sinais, a Fase 00 tem $(138 + 12)$ execuções para a voz e $(46 + 12)$ para o EEG, isto é, $150 + 58 = 208$ execuções de ranking. O vetor fundido não é varrido aqui: ele é montado na Fase 01, a partir dos vencedores de cada sinal.

			\subsubsection{Fase 01 --- autenticação com o vetor vencedor}
				\label{sec:fase01}

				\par A Fase 01 usa \emph{apenas} a combinação vencedora da Fase 00 --- o eixo do extrator não é mais variado, pois a engenharia paraconsistente já o fixou. Variam-se quatro eixos, detalhados adiante: a \textbf{fonte} do sinal, o \textbf{modo de texto}, o \textbf{esquema de validação cruzada} e a \textbf{padronização das características}. São $4 \times 2 \times 2 \times 2 = 32$ execuções, todas com o classificador \ac{DSNN}. O \autoref{lst:e05Phase01} mostra a montagem da fonte e o laço de autenticação; a \autoref{fig:e05Phase01Flow} resume o fluxo.

				\paragraph{Arquitetura do classificador DSNN.} O classificador é uma rede de pulsos residual profunda, construída a partir da especificação declarativa \texttt{["linear:128:relu", "residual:2", "linear:N:identity"]} do perfil, onde $N$ é o número de locutores. Concretamente, ela é composta por:
				\begin{itemize}
					\item um \textbf{estágio de entrada} \texttt{Linear}($d \to 128$) seguido de um neurônio \ac{LIF}, onde $d$ é a dimensão do vetor de características vindo da Fase 00;
					\item \textbf{dois blocos residuais} idênticos, cada um \texttt{Linear}($128 \to 128$) seguido de \ac{LIF}, envoltos por uma conexão de atalho identidade, $h_{\text{saída}} = \text{bloco}(h) + h$. Os atalhos não têm parâmetros e existem para que o gradiente alcance as camadas iniciais sem se dissipar através da não-linearidade de disparo --- é o mesmo desenho de \ac{SNN} residual profunda de \cite{zheng2021goingdeeper}, para o qual a \ac{TDBN} foi originalmente proposta. Os estágios de entrada e de saída não recebem atalho porque suas larguras diferem de 128, o que impediria a soma elemento a elemento;
					\item um \textbf{estágio de saída} \texttt{Linear}($128 \to N$), cuja saída é a média dos logits sobre os $T$ passos de tempo.
				\end{itemize}
				O treino usa retropropagação através do tempo (\ac{BPTT}) com gradiente substituto, e a regularização da taxa de disparo descrita adiante mantém cada camada pulsante fora dos regimes de neurônio morto e saturado.

				\paragraph{Protocolo de verificação: como EER e AUC são calculados.} Esta é a parte do protocolo que a definição textual de \ac{EER} não revela, e sem a qual os números da \autoref{chap:testsResults} não são interpretáveis. A camada de saída tem uma unidade por locutor de \emph{treino}, mas ela \textbf{não} é usada como decisão de classe --- sob o protocolo de verificação os locutores de teste nunca aparecem no treino, de modo que nenhuma dessas unidades corresponde ao indivíduo sendo verificado. O vetor de logits é usado, em vez disso, como um \textbf{vetor de representação} (\textit{embedding}) do sinal. A partir dele:
				\begin{enumerate}
					\item cada par de amostras de teste forma uma \textbf{tentativa} (\textit{trial}), rotulada \textit{genuína} quando as duas amostras pertencem ao mesmo locutor e \textit{impostora} caso contrário;
					\item a pontuação de cada tentativa é a \textbf{similaridade de cosseno} entre os dois vetores de representação;
					\item varre-se um limiar sobre essas pontuações; o \ac{EER} é o ponto em que a taxa de falsa aceitação iguala a de falsa rejeição;
					\item a \ac{AUC} é estimada pela estatística de Wilcoxon--Mann--Whitney, $\text{AUC} = P(\text{pontuação genuína} > \text{pontuação impostora})$, que equivale à área sob a curva \ac{ROC} das duas distribuições de pontuação.
				\end{enumerate}
				É por isso que as métricas de conjunto fechado são \texttt{NaN} nos resultados: elas exigiriam que a identidade de teste estivesse entre as classes treinadas, o que este protocolo deliberadamente proíbe.

				\paragraph{Os quatro eixos experimentais.} Cada eixo é binário ou quaternário, e o experimento é o produto cartesiano completo:
				\begin{itemize}
					\item \textbf{Fonte} (4 níveis): apenas voz, apenas EEG, fundida-precoce e fundida-tardia --- os três modos da \autoref{sec:estruturaDaEstrategiaProposta}, com o modo fundido desdobrado nas duas variantes de instante de fusão;
					\item \textbf{Modo de texto} (2 níveis): \textit{dependente}, em que o mesmo enunciado aparece no treino e no teste, e \textit{independente}, em que os enunciados de teste podem ser quaisquer outros. É o eixo da Questão~4 (\autoref{chap:intro}), e mede o quanto do desempenho vem da identidade do locutor e não do conteúdo falado;
					\item \textbf{Esquema de validação cruzada} (2 níveis). Em ambos os esquemas as partições são \emph{disjuntas por locutor} (\textit{group $k$-fold}, $k=5$), o que é o que torna o protocolo de verificação válido. Na \textbf{plana}, cada partição é treinada e pontuada uma vez, sem conjunto de validação e sem parada antecipada. Na \textbf{aninhada}, um laço interno --- também $k=5$, com semente distinta --- separa um conjunto de validação usado para selecionar o melhor modelo por parada antecipada, e só então o modelo escolhido é pontuado uma única vez na partição externa. A aninhada é o estimador não enviesado quando há seleção de modelo; a plana é mais barata e não descarta dados para validação. Os dois são comparados justamente porque a diferença entre eles mede o custo dessa seleção;
					\item \textbf{Padronização das características} (2 níveis): \textit{z-score} por característica ligado ou desligado. Quando ligado, média e desvio padrão são estimados \emph{apenas} sobre a partição de treino de cada \textit{fold} e então aplicados à de teste, evitando vazamento de informação. É uma ablação do efeito da normalização.
				\end{itemize}

				\paragraph{Protocolo de treino e de repetição.} Todas as 32 execuções compartilham os mesmos hiperparâmetros: otimizador Adam com taxa de aprendizado $10^{-3}$, 50 épocas, lotes de 32 amostras, e paciência de 10 épocas para a parada antecipada (ativa apenas na validação cruzada aninhada, já que a plana não reserva conjunto de validação). Cada configuração é executada \textbf{três vezes}, com sementes distintas derivadas da semente base $42$; os valores reportados no \autoref{chap:testsResults} são as médias sobre essas três repetições, e o desvio padrão ali tabelado é calculado entre elas. Como cada repetição percorre $k=5$ partições, cada configuração contribui com 15 medidas de \ac{EER}.

				\begin{lstlisting}[language=C, caption={Fase 01 --- autenticação com o vencedor da Fase 00. O extrator nao e mais variado: usa-se apenas a combinacao escolhida por sinal.}, label={lst:e05Phase01}]
/* v_voz, v_eeg: vencedores da Fase 00 (um por sinal)            */
/* fonte : VOZ | EEG | FUNDIDA_PRECOCE | FUNDIDA_TARDIA          */
/* Preenche X (vetores de caracteristicas) e devolve a dimensao. */

int monta_fonte(Amostra a, int fonte,
                Combinacao v_voz, Combinacao v_eeg, double X[])
{
    if (fonte == VOZ)
        return extrai(a.audio, v_voz, X);
    if (fonte == EEG)
        return extrai(a.eeg,   v_eeg, X);

    if (fonte == FUNDIDA_PRECOCE) {         /* funde sinal bruto */
        double s[MAX_AMOSTRAS];
        int n = concatena(a.audio, a.eeg, s);
        return extrai(s, v_voz, X);         /* uma extracao so   */
    }
    /* FUNDIDA_TARDIA: extrai cada sinal, concatena os vetores   */
    double xv[MAX_CARACT], xe[MAX_CARACT];
    int mv = extrai(a.audio, v_voz, xv);
    int me = extrai(a.eeg,   v_eeg, xe);
    return concatena_vetores(xv, mv, xe, me, X);
}

/* Loop de autenticacao: fonte X modo de texto X esquema de CV.  */
Metricas autentica(Amostra base[], int fonte,
                   int modo_texto, int cv_aninhada,
                   Combinacao v_voz, Combinacao v_eeg)
{
    double X[MAX_AMOSTRAS][MAX_CARACT];
    for (int i = 0; i < N_AMOSTRAS; i++)
        monta_fonte(base[i], fonte, v_voz, v_eeg, X[i]);

    /* grupos disjuntos por locutor (protocolo de verificacao)   */
    Folds f = divide_por_locutor(base, modo_texto, cv_aninhada);

    Metricas m = treina_e_avalia_dsnn(X, f);  /* LIF + BPTT       */
    return m;                                 /* EER e AUC        */
}
\end{lstlisting}


				\begin{figure}[h]
					\centering
					\caption[Fluxograma da Fase 01]{Autenticação DSNN na Fase 01}
					\scalebox{0.82}{\begin{tikzpicture}[
    startstop/.style={
        rectangle, rounded corners=4pt, draw, fill=flowStartEndColor,
        text=black, minimum width=3.6cm, minimum height=0.7cm,
        font=\footnotesize\sffamily\bfseries, align=center
    },
    proc/.style={
        rectangle, rounded corners, draw, fill=flowProcessColor,
        minimum width=5.6cm, minimum height=0.7cm,
        font=\footnotesize\sffamily\bfseries, align=center
    },
    src/.style={
        rectangle, rounded corners, draw, fill=flowProcessColor!55,
        minimum width=5.6cm, minimum height=0.7cm,
        font=\scriptsize\sffamily\bfseries, align=center
    },
    arr/.style={->, >=stealth, draw=flowArrowColor, line width=1.0mm},
    node distance=0.5cm
]

\node[startstop] (win) {Vetores vencedores da Fase 00 (voz, EEG)};
\node[src, below=of win] (build)
    {montar a fonte:\\ voz $\mid$ EEG $\mid$ fundida-precoce $\mid$ fundida-tardia};
\node[proc, below=of build] (encode)
    {codificar por injeção de corrente ($T$ passos)};
\node[proc, below=of encode] (dsnn)
    {classificador DSNN residual (LIF + BPTT)\\ leitura $=$ taxa média de disparo};
\node[proc, below=of dsnn] (cv)
    {validação cruzada por locutor (aninhada ou plana),\\ modo de texto dependente ou independente};
\node[startstop, below=of cv] (metrics) {EER e AUC (protocolo de verificação)};

\draw[arr] (win)    -- (build);
\draw[arr] (build)  -- (encode);
\draw[arr] (encode) -- (dsnn);
\draw[arr] (dsnn)   -- (cv);
\draw[arr] (cv)     -- (metrics);
\end{tikzpicture}
}
					\legend{Fonte: Elaborado pelo autor, 2026.}
					\label{fig:e05Phase01Flow}
				\end{figure}

			\subsubsection{Organização dos perfis e transição entre fases}

				\par Cada execução é descrita por um \textit{perfil} em formato \ac{JSON}. Os perfis ficam separados por fase: 208 perfis em \texttt{profiles/phase00/} (um por combinação varrida) e 32 em \texttt{profiles/phase01/} (fonte $\times$ texto $\times$ \ac{CV} $\times$ padronização). O vencedor da Fase 00 só é conhecido após a execução; por isso os perfis da Fase 01 nascem com um extrator provisório, substituído pelo vencedor antes de rodar.

				\par Dois \textit{scripts} automatizam a transição, evitando edição manual (\autoref{fig:e05PipelineTwoPhase}):
				\begin{itemize}
					\item \texttt{thesis\_phase00\_rank.py}: lê os resultados paraconsistentes das 208 execuções, calcula a média de $D_{\text{penalizado}}$ sobre as três repetições, escreve o vencedor de cada sinal em \texttt{winners.json} e sinaliza explicitamente eventuais empates exatos, para que uma escolha feita por desempate de ordenação não seja confundida com um resultado de mérito;
					\item \texttt{thesis\_apply\_winner.py}: injeta cada vencedor nos 32 perfis da Fase 01 --- voz recebe o vencedor da voz, EEG o do EEG, e as fontes fundidas recebem, por padrão, o vencedor da voz (sinal dominante em número de amostras).
				\end{itemize}

		\subsection{Métricas}
		
			\par A autenticação segue o protocolo de \textbf{verificação} (\textit{verification}), no qual os locutores de teste não aparecem no treino: os grupos de validação cruzada são disjuntos por locutor. Nesse protocolo, as métricas primárias são a \textbf{\ac{EER}} (o ponto em que a taxa de falsa aceitação iguala a de falsa rejeição) e a \textbf{área sob a curva \ac{ROC}} (\ac{AUC}). As métricas de classificação de conjunto fechado (acurácia, F1, precisão, recall, especificidade) só se aplicam quando as classes de teste também estão no treino; sob o protocolo de verificação elas não são reportadas --- nos arquivos de resultado da Fase~01 elas aparecem literalmente como \texttt{NaN}, e não como zero, justamente para que a ausência de valor não seja confundida com um desempenho nulo. As seguintes métricas são adotadas, garantindo a comparação com outros estudos:

			\begin{itemize}
				\item \textbf{Taxa de Erro Igual (EER)} --- primária, protocolo de verificação;
				\item \textbf{área sob a Curva ROC (AUC)} --- primária, protocolo de verificação;
				\item acurácia, pontuação F1, precisão, recall, sensitividade, especificidade --- reportadas apenas em avaliações de classificação de conjunto fechado;
				\item Erro Médio Quadrático (EQM) --- usado na reconstrução dos \textit{autoencoders}, não na classificação.
			\end{itemize}
			
		\subsection{Desenvolvimento e Validação de Algoritmos}
		
			\par Quanto aos algoritmos de classificação, empregam-se redes neurais de pulso em uma arquitetura residual (a \ac{DSNN} detalhada na \autoref{sec:fase01}), com o objetivo de minimizar o gasto de energia, memória e processamento, adequando-se a dispositivos com baixo poder computacional para garantir portabilidade e eficiência energética.

			\par Há, no entanto, um obstáculo prático: o número de locutores e de amostras por locutor é pequeno. Com poucos dados, a rede tende a ``decorar'' os exemplos de treino em vez de aprender padrões gerais --- o chamado \textit{overfitting}. Para combatê-lo, usam-se duas técnicas de regularização complementares, descritas em detalhe no \autoref{chap:regularizacao}.

			\par A primeira é o \textbf{decaimento de peso L2}, que penaliza pesos muito grandes e, com isso, favorece soluções mais simples e que generalizam melhor. Nas redes de pulso, atua apenas sobre as matrizes de peso; os parâmetros biofísicos dos neurônios de pulso ($R$, $C$ e $V_{th}$) são preservados, pois reduzi-los descaracterizaria o neurônio.

			\par A segunda é a \textbf{regularização da taxa de disparo}, exclusiva das redes de pulso. Ela mantém a frequência média de disparos de cada camada dentro de uma faixa saudável, evitando dois extremos: neurônios mortos, que param de disparar e interrompem o aprendizado, e neurônios saturados, que disparam o tempo todo e perdem a capacidade de distinguir entradas. Como bônus, ao incentivar disparos comedidos, essa técnica reforça a própria razão de ser das redes de pulso: a eficiência energética.
			
%	\section{Cronograma}
%		\label{sec:cronograma}
%	
%		\begin{table}[h!]
%			\centering
%			\caption{Datas dos trabalhos a serem realizados}
%			\tableZebraStripe
%			\begin{tabular}{@{}Y{10cm}cc@{}}
%				\tableHeaderRow \textbf{Etapa} & \textbf{Início} & \textbf{Fim} \\
%				Programação do sistema de coleta de dados & 01/09/2024 & 01/12/2024 \\
%				Extração e seleção de características & 01/02/2025 & 01/12/2025 \\
%				Desenvolvimento de classificadores & 01/12/2025 & 01/12/2026 \\
%				Validação experimental e comparação & 01/03/2026 & 01/02/2027 \\
%				Divulgação e publicação dos resultados & 01/02/2027 & 01/04/2027 \\
%				Redação da dissertação e preparação para defesa & 01/09/2024 & 01/03/2027 \\
%				\bottomrule
%			\end{tabular}
%			\label{tab:cronograma}
%		\end{table}
%		
%		\begin{table}[h!]
%			\centering
%			\caption{Gráfico de Gantt dos trabalhos a serem realizados}
%			\begin{ganttchart}[
%				hgrid, % Adiciona linhas horizontais de grade para melhorar a legibilidade do gráfico
%				vgrid, % Adiciona linhas verticais de grade
%				x unit=.455cm, % Define a largura de cada unidade de tempo (neste caso, um mês) em centímetros
%				y unit chart=0.8cm, % Define a altura de cada linha de tarefa no gráfico em centímetros
%				bar/.append style={fill=ganttBarColor, draw=ganttBarColor}, % Estilo padrão para as barras (paleta colorblind-safe, visualStyle.tex)
%				bar label font=\scriptsize\bfseries, % Define o estilo da fonte para os rótulos das barras (pequeno e em negrito)
%				bar label node/.append style={text=black, anchor=west, xshift=0.2cm}, % Alinha o texto da barra ao início da barra com um pequeno deslocamento horizontal
%				group label node/.append style={anchor=east, font=\small}, % Define o estilo dos rótulos de grupo com ânfora à esquerda e fonte pequena
%				time slot format=isodate, % Define o formato de data para os slots de tempo (ISO: ano-mês-dia)
%				milestone/.append style={fill=ganttMilestoneColor}, % Estilo para marcos
%				title/.append style={fill=tableHeaderColor, draw=black}, % Estilo do título: preenchimento azul e borda preta
%				title label font=\bfseries\footnotesize, % Define o estilo da fonte para os rótulos do título (pequeno e em negrito)
%				title label anchor/.style={below=0ex}, % Ajusta a âncora do rótulo do título para baixo
%				title/.style={fill=ganttTitleColor, draw=none}, % Estilo do título: preenchimento azul claro e sem borda
%				title label font=\scriptsize\bfseries, % Define o estilo da fonte para os rótulos do título (pequeno e em negrito)
%				milestone label font=\small\bfseries, % Define o estilo da fonte para os rótulos de marcos (pequeno e em negrito)
%				bar height=0.5, % Define a altura das barras
%				time slot unit=month, % Define a unidade de tempo como mês
%				y unit title=0.9cm % Define a altura da linha do título em centímetros
%				]{2024-08-10}{2027-05-01} % Define o intervalo de datas para o gráfico Gantt
%				
%				\gantttitlecalendar{year, month=number} \\ % Adiciona um título de calendário ao gráfico, mostrando ano e número do mês
%				
%				% Define as barras de tarefas com seus rótulos e datas de início e fim
%				\ganttbar[
%				bar label node/.append style={xshift=6pt} % Ajusta o texto da barra para começar no início da barra
%				]{Programação do sistema de coleta de dados}{2024-09-01}{2024-12-01} \\
%				\ganttbar[
%				bar label node/.append style={xshift=70pt} % Ajusta o texto da barra para começar no início da barra
%				]{Extração e seleção de características}{2025-02-01}{2025-12-01} \\
%				\ganttbar[
%				bar label node/.append style={xshift=200pt} % Ajusta o texto da barra para começar no início da barra
%				]{Desenvolvimento de classificadores}{2025-12-01}{2026-12-01} \\
%				\ganttbar[
%				bar label node/.append style={xshift=240pt} % Ajusta o texto da barra para começar no início da barra
%				]{Validação experimental e comparação}{2026-03-01}{2027-02-01} \\
%				\ganttbar[
%				bar label node/.append style={xshift=230pt} % Ajusta o texto da barra para começar no início da barra
%				]{Divulgação e publicação dos resultados}{2027-02-01}{2027-04-01} \\
%				\ganttbar[
%				bar/.append style={fill=ganttHighlightColor}, % Barra em destaque (paleta colorblind-safe, visualStyle.tex)
%				bar label node/.append style={xshift=4pt} % Ajusta o texto da barra para começar no início da barra
%				]{Redação da dissertação e preparação para defesa}{2024-09-01}{2027-03-01}
%			\end{ganttchart}
%			\label{tab:cronogramaGantt}
%		\end{table}
%
%\section{Resultados Esperados}
%
%	\par Esperava-se desenvolver algoritmos de extração de características e classificação para voz e eletroencefalograma (EEG) que aprimorassem a acurácia na identificação de indivíduos; avaliar se as alegações sobre o consumo energético e a eficácia das redes neurais de pulso na classificação de informações temporais se confirmam para o problema estudado; e comparar o desempenho das técnicas de extração de características manuais com as automatizadas. O \autoref{chap:testsResults} confronta cada uma dessas expectativas com o que foi de fato medido, e o \autoref{chap:conclusions} discute em que medida elas se sustentaram.


%\section{Procedimentos}
%\subsection{Tratamento do sinal}
%\subsection{Procedimento 01}